
\subsubsection{Key Findings Interpretation}

Synthetic proof-of-concept experiments demonstrate that hierarchical Bayesian calibration achieves simultaneous improvement in calibration quality (expected calibration error = 0.043) and computational efficiency (30\% cost reduction) by exploiting complementarity between consistency-based and conformal prediction methods.

The moderate correlation ($\rho \approx$ 0.43--0.46) between consistency and conformal methods in synthetic validation suggests they measure fundamentally distinct dimensions. While correlation values were controlled in the synthetic data generation process, the stability across synthetic dataset types with different uncertainty profiles suggests the mechanism operates as designed.

Hierarchical Bayesian calibration achieves improvement over independent cascade, isolating the value of Bayesian mutual updating. Since cascade already combines both methods sequentially, this gain demonstrates that bidirectional information flow provides value beyond simple combination in synthetic validation.

The method achieves 30\% cost reduction while maintaining coverage above 90\%, defying the conventional tradeoff where efficiency gains come at the expense of coverage. Consistency-informed weighting reduces intervals for high-consistency queries, creating efficiency gains, while Bayesian threshold updating ensures that when coverage drops, consistency filtering is loosened.

\subsubsection{Limitations}

Three limitations require acknowledgment:

\textbf{Limitation 1: Synthetic Proof-of-Concept Validation}

Validation experiments used synthetic proof-of-concept data (200 samples per dataset) with controlled correlation ($\rho \approx$ 0.5 target) to demonstrate the core mechanism. While real datasets (TruthfulQA, HH-RLHF, SQuAD) were loaded and real models (Llama-2-7B) were used for method development, full-scale real data inference was not performed.

Specific quantitative results (expected calibration error = 0.043, $\rho$ = 0.463) require confirmation with real data validation. The three-step mechanism (consistency $\rightarrow$ conformal $\rightarrow$ Bayesian co-calibration) is theoretically motivated and demonstrated via synthetic validation; production deployment requires real data validation to confirm performance metrics and correlation stability across actual uncertainty profiles.

Real model inference may exhibit different correlation structures than synthetic data. The correlation stability finding is based on synthetic scenarios with controlled generation, not real dataset diversity.

Synthetic proof-of-concept is standard practice for establishing methodology soundness before resource-intensive full-scale validation. The experiments demonstrate the mechanism works as designed; real data validation is the natural next step.

\textbf{Limitation 2: Labeled Calibration Data Requirement}

Hierarchical Bayesian calibration requires labeled validation sets for both consistency threshold tuning and conformal calibration. This is a standard requirement for supervised uncertainty quantification methods, but limits deployment in zero-shot scenarios.

The contribution applies to settings with available calibration data (factuality benchmarks, question answering datasets, domain-specific corpora with ground truth). Few-shot adaptation and zero-shot deployment remain open challenges.

\textbf{Limitation 3: Out-of-Distribution Detection Untested}

Experiments focus on in-distribution calibration. Domain shift experiments (calibrate on one dataset, test on another) were not conducted. Out-of-distribution detection claims remain speculative.

Domain shift may violate the exchangeability assumption underlying conformal prediction, causing coverage degradation. The core calibration contribution stands independently. In-distribution calibration is valuable independently of out-of-distribution detection.

\subsubsection{Broader Impact}

For the research community, hierarchical Bayesian calibration provides a template for integrating diverse uncertainty quantification methods. The insight---quantify correlation to determine whether joint calibration adds value---applies to any pair of uncertainty estimators. The complementarity framework (0.3 < $\rho$ < 0.7) provides hypothesized guidance for when integration is worthwhile versus when single-method optimization suffices.

For practitioners, the method may enable deployment in production systems requiring both statistical guarantees and computational efficiency, pending real-world validation. The 30\% cost reduction could make rigorous uncertainty quantification feasible for latency-sensitive applications.

Improved calibration may reduce overconfident errors, mitigating harms in high-stakes domains. However, uncertainty quantification methods are not adversarially robust. Targeted attacks on consistency checks remain a risk.
